\chapter{A First Session}
\label{ch:tutorial}

This chapter builds a tiny rule base at the shell prompt, one step at a
time. Start \fn{./morendo -shell} and type along; every transcript below is
real output. Chapter~\ref{ch:language} then describes each construct in
full.

\section{Templates and facts}

Every fact in \Morendo\ has a template that names its slots. Define one
with \fn{deftemplate}, then create facts with \fn{assert}:

\begin{lstlisting}[style=shell]
Morendo> (deftemplate person (slot name) (slot age))
true
Morendo> (assert (person (name "ann") (age 34)))
2
Morendo> (assert (person (name "bob") (age 12)))
3
Morendo> (facts)
f-1 (_initialFact)
f-2 (person (name "ann") (age 34) )
f-3 (person (name "bob") (age 12) )
for a total of 3
\end{lstlisting}

\fn{assert} returns the id of the new fact. Fact~1 is the
\emph{initial fact}, which the engine creates itself; rules whose first
pattern is negated or that have no patterns at all are matched against it.
Slots you leave out get their default value, which is \fn{nil} unless the
template says otherwise. \fn{(clear)} empties the engine, templates
included, so the template can be defined again with a default:

\begin{lstlisting}[style=shell]
Morendo> (clear)
true
Morendo> (deftemplate person (slot name) (slot age (type INTEGER) (default 0)))
true
Morendo> (assert (person (name "ann") (age 34)))
2
Morendo> (assert (person (name "bob") (age 12)))
3
Morendo> (assert (person (name "cy")))
4
Morendo> (facts)
f-1 (_initialFact)
f-2 (person (name "ann") (age 34) )
f-3 (person (name "bob") (age 12) )
f-4 (person (name "cy") (age 0) )
for a total of 4
\end{lstlisting}

Slot values are strings in double quotes, integers, floating-point numbers,
the booleans \fn{TRUE} and \fn{FALSE}, or \fn{nil}. Symbols without quotes
are accepted where CLIPS accepts them, but strings are the usual choice for
text.

\section{A rule}

A rule has a name, an optional comment string, a left-hand side of
patterns, the arrow \fn{=>}, and a right-hand side of actions. This one
prints every person who is at least 18:

\begin{lstlisting}
(defrule adult "people who may vote"
  (person (name ?n) (age ?a&:(>= ?a 18)))
=>
  (printout t ?n " is an adult" crlf))
\end{lstlisting}

The pattern \fn{(person ...)} matches facts of the \fn{person} template.
\fn{?n} and \fn{?a} are variables; a variable in a slot binds to the slot's
value. \fn{\&:} attaches a \emph{predicate constraint}: the pattern only
matches when the expression after the colon is true. \fn{printout t} writes
to the standard output; \fn{crlf} ends the line.

Type the rule into the shell (the prompt shows \fn{...} while it waits for
the closing parentheses) and list the rules:

\begin{lstlisting}[style=shell]
Morendo> (rules)
adult "people who may vote" salience:100 version: no-agenda:false temporal-activation:false cost-value:0
for a total of 1
\end{lstlisting}

Nothing has been printed yet. Matching happened when the rule was
compiled, but actions run only when you fire the agenda:

\begin{lstlisting}[style=shell]
Morendo> (watch rules)
Morendo> (fire)
==> fire: Activation: adult, id-2 AggrTime--1
ann is an adult
1
Morendo> (fire)
0
Morendo> (unwatch rules)
\end{lstlisting}

\fn{(watch rules)} makes the engine announce each rule as it fires, naming
the facts that activated it (\fn{id-2} is Ann's fact), and \fn{unwatch}
turns the trace off again. \fn{fire} returns the
number of rules fired; the second call finds an empty agenda. A rule fires
once per combination of facts that matches it; it does not fire again for
the same facts unless they are retracted and asserted anew, or modified.

\section{Retracting and modifying facts}

\fn{retract} takes a fact id or a fact bound to a variable; \fn{modify}
takes a bound fact and new slot values. In the shell, bind the value of
\fn{assert} to keep a handle on the fact:

\begin{lstlisting}[style=shell]
Morendo> (watch facts)
Morendo> (bind ?f (assert (person (name "dee") (age 40))))
==> f-5 (person (name "dee") (age 40) )
true
Morendo> (modify ?f (age 41))
<== f-5 (person (name "dee") (age 40) )
==> f-5 (person (name "dee") (age 41) )
true
Morendo> (retract ?f)
<== f-5 (person (name "dee") (age 41) )
true
Morendo> (unwatch facts)
\end{lstlisting}

With \fn{(watch facts)} on, \fn{==>} marks an assertion and \fn{<==} a
retraction. A modification is a retraction followed by an assertion of the
changed fact under the same id, and every rule that depends on the fact is
re-evaluated. \fn{(retract 2)} with a literal id works as well, and a
slot value may be computed in place: \fn{(modify ?f (age (+ 40 1)))}
does what the transcript did without binding the value first.

\section{Salience and the agenda}
\label{sec:tutorial-salience}

When several activations are waiting, \emph{salience} decides which fires
first. Every rule has salience 100 unless it declares otherwise; higher
fires first. Define two rules, start from an empty working memory with
\fn{(reset)}, and assert two facts:

\begin{lstlisting}
(defrule greet (person (name ?n)) => (printout t "hello " ?n crlf))
(defrule low (declare (salience 50))
  (person (name ?n)) => (printout t "hi " ?n crlf))
\end{lstlisting}

\begin{lstlisting}[style=shell]
Morendo> (reset)
Morendo> (assert (person (name "ed") (age 1)))
2
Morendo> (watch activations)
Morendo> (assert (person (name "flo") (age 2)))
=> Activation: greet, id-3 AggrTime--1
=> Activation: low, id-3 AggrTime--1
3
Morendo> (unwatch activations)
Morendo> (agenda)
100 greet: f-3
100 greet: f-2
50 low: f-3
50 low: f-2
for a total of 4
Morendo> (fire 1)
hello flo
1
Morendo> (agenda)
100 greet: f-2
50 low: f-3
50 low: f-2
for a total of 3
Morendo> (fire)
hello ed
hi flo
hi ed
3
\end{lstlisting}

\fn{(watch activations)} shows activations being added to the agenda as
the facts are asserted, and \fn{(agenda)} shows the queue in firing
order. All of \fn{greet}'s activations (salience 100) come before
\fn{low}'s (salience 50). Within one salience level the default strategy
is depth-first: the most recently added activation fires first, which is
why \fn{flo} precedes \fn{ed}. \fn{(fire 1)} fires only the first
activation; \fn{(fire)} then runs the rest. \fn{set-strategy} selects
breadth or recency ordering, and \fn{(refresh rule)} puts a rule back on
the agenda for matches it has already fired on (Section~\ref{sec:agenda}).

\section{Rules and facts in either order}
\label{sec:rules-before-facts}

Matching happens at assertion time, but a rule defined later is matched
against the facts that already exist: when the rule's nodes are added to
the network the engine replays the facts held by the nodes they hang off.
Both rules below fire for the fact asserted before them:

\begin{lstlisting}[style=shell]
Morendo> (clear)
true
Morendo> (deftemplate person (slot name) (slot age))
true
Morendo> (assert (person (name "cy") (age 0)))
2
Morendo> (defrule lit (person (name "cy")) => (printout t "literal matched" crlf))
true
Morendo> (defrule var (person (name ?n)) => (printout t "variable matched " ?n crlf))
true
Morendo> (fire)
literal matched
variable matched cy
2
\end{lstlisting}

Templates, however, must exist before the rules and facts that use them:
a rule naming an unknown template is rejected, and so is an \fn{assert}.
Rule files therefore put their \fn{deftemplate}s first, and by convention
their \fn{defrule}s next and the facts last, as every file in
\file{samples/} does.

\section{Functions, variables and control}

Outside rules the shell is a calculator. Function calls are prefix
expressions; \fn{bind} sets a variable that stays defined for the rest of
the session:

\begin{lstlisting}[style=shell]
Morendo> (* 2.5 4)
10.0
Morendo> (bind ?x (* 3 4))
true
Morendo> (printout t "x is " ?x crlf)
x is 12
Morendo> (if (> ?x 10) then (printout t "big" crlf) else (printout t "small" crlf))
big
Morendo> (now)
2026-09-12T16:37:33.845210878Z
\end{lstlisting}

\fn{deffunction} defines a function in CLIPS itself, and \fn{defglobal} a
global variable, written with asterisks around its name:

\begin{lstlisting}[style=shell]
Morendo> (deffunction twice (?n) (* 2 ?n))
true
Morendo> (twice 21)
42
Morendo> (defglobal ?*limit* 18)
18
Morendo> (printout t ?*limit* crlf)
18
\end{lstlisting}

The CLIPS form \fn{(defglobal ?*limit* = 18)} is accepted as well.
For output with a fixed layout, \fn{format} takes a C-style format
string: \fn{(format t "\%-6s \%5.2f\%n" "x" ?x)}
(Section~\ref{sec:routers}).

Loops are functions too: \fn{while} repeats while a test holds,
\fn{loop-for-count} counts, and \fn{foreach} walks a list. A
\fn{deffunction} body may hold several expressions and \fn{return}
early:

\begin{lstlisting}[style=shell]
Morendo> (loop-for-count (?k 1 3) (printout t "k=" ?k crlf))
k=1
k=2
k=3
false
Morendo> (foreach ?e (create$ a b c) (printout t ?e-index ":" ?e crlf))
1:a
2:b
3:c
false
Morendo> (deffunction size (?a ?b)
  (bind ?c (+ ?a ?b))
  (if (> ?c 10) then (return "big"))
  "small")
true
Morendo> (size 8 9)
big
\end{lstlisting}

Section~\ref{sec:loops} has the full set, \fn{progn}, \fn{progn\$} and
\fn{break} included.

\section{Rule files}

Rather than typing at the prompt, put constructs in a file and load it
with \fn{batch}. The file \file{only\_1.clp} in \file{samples/only/}
is typical: two templates, one rule, then the facts. \fn{(batch file)} evaluates every
expression in the file in order and returns \fn{true}; anything the file
prints appears as it is loaded.

\begin{lstlisting}[style=shell]
Morendo> (batch samples/only/only_1.clp)
true
Morendo> (fire)
spiderman save the day
1
\end{lstlisting}

Two other functions load data: \fn{load-facts} reads a file of fact
expressions (no templates or rules) and asserts them, which is how the
Manners benchmark loads its guest lists; \fn{load-graph} reads nodes and
edges for graph queries (Chapter~\ref{ch:queries}). A file can also
declare its starting facts in a \fn{deffacts} and end with \fn{(reset)},
which retracts whatever is in working memory and asserts them; this is
the CLIPS habit, and it makes a rule file rerunnable
(Section~\ref{sec:deffacts}).

\section{Inspecting the engine}

\begin{center}
\begin{tabularx}{\textwidth}{lL}
\toprule
\fn{(facts)} & list working memory \\
\fn{(agenda)} & list the activations waiting to fire, in order \\
\fn{(rules)} & list the rules with their salience and properties \\
\fn{(list-deffacts)} & list the deffacts \\
\fn{(templates)} & list the templates, including the built-in \fn{Node}, \fn{Edge} and \fn{Graph} \\
\fn{(ppdefrule name)} & pretty-print a rule as the engine compiled it, with every property and its topology cost \\
\fn{(ppdeftemplate name)} & pretty-print a template \\
\fn{(watch x)}, \fn{(unwatch x)} & trace \fn{facts}, \fn{rules}, \fn{activations} or \fn{all} \\
\fn{(profile x)}, \fn{(unprofile x)} & time \fn{assert-fact}, \fn{retract-fact}, \fn{fire}, \fn{add-activation}, \fn{remove-activation} or \fn{all} \\
\fn{(clear)} & remove every fact, rule and template \\
\fn{(exit)} & close the engine and leave the shell \\
\bottomrule
\end{tabularx}
\end{center}

\fn{ppdefrule} is worth a look after writing a rule: it shows the full
\fn{declare} block with every property the rule has, explicitly or by
default:

\begin{lstlisting}[style=shell]
Morendo> (ppdefrule var)
(defrule var
  (declare (salience 100) (remember-alpha true) (chaining-direction forward)
           (no-agenda false) (temporal-activation false) )
  (person
    (name ?n)
  )
=>
  (printout t "variable matched " ?n crlf)
)
;; topology-cost: 0
\end{lstlisting}

The declaration is printed in the form the parser accepts, so the output
can be pasted back into a rule file.
